\documentclass[11pt,a4paper]{article}
\usepackage[margin=25mm]{geometry}
\usepackage{fontspec}
\setmainfont{TeX Gyre Pagella}
\usepackage{amsmath,amssymb,mathtools}
\usepackage{graphicx}
\usepackage{xcolor}
\usepackage[unicode,colorlinks=true,linkcolor=blue,urlcolor=blue]{hyperref}
\usepackage{bookmark}
\usepackage{enumitem}
\setlist{itemsep=3pt,topsep=5pt}
\setlength{\emergencystretch}{3em}
\newcommand{\tightlist}{\setlength{\itemsep}{0pt}\setlength{\parskip}{0pt}}
\newcommand{\passthrough}[1]{#1}
\newcommand{\pandocbounded}[1]{#1}
\makeatletter
\def\maxwidth{\ifdim\Gin@nat@width>\linewidth\linewidth\else\Gin@nat@width\fi}
\def\maxheight{\ifdim\Gin@nat@height>0.75\textheight0.75\textheight\else\Gin@nat@height\fi}
\makeatother
\setkeys{Gin}{width=\maxwidth,height=\maxheight,keepaspectratio}
\hypersetup{pdfauthor={OpenAI GPT-6.1 Sol, Codex, Ultra effort}}
\begin{document}
\section{Weil groups and one-dimensional
representations}\label{weil-groups-and-one-dimensional-representations}

\emph{Written by OpenAI GPT-6.1 Sol in Codex, Ultra effort, October
2026. Self-checked by the writing AI; no independent review is claimed.
Public domain (CC0).}

Local reciprocity embeds \(K^\times\) densely in the abelianized
absolute Galois group. Its valuation coordinate is an integer, whereas
the Galois group has a profinite valuation coordinate. The Weil group
restores that integer coordinate and its discrete topology. Reciprocity
then becomes a homeomorphism, and multiplicative characters become
one-dimensional Weil representations.

We use \href{../explicit-local-reciprocity-and-existence.html}{Explicit
local reciprocity and the existence theorem}, especially Theorem 9.4,
and the unit-filtration theorem in
\href{../abelian-ramification-conductors-and-hasse-arf.html}{Abelian
ramification, conductors and Hasse--Arf}. The analytic character factors
come from \emph{Tate's local theory at the finite places} and
\emph{Tate's local theory at the infinite places} in the course
\emph{Adèles and L-functions}, lessons 7 and 8. We supply the additional
Fourier argument needed in positive characteristic. Sections 1--4
concern a nonarchimedean local field of either characteristic. Section 5
treats the real and complex fields; section 6 translates Frobenius
conventions, and section 7 connects the local construction to global
reciprocity.

\textbf{Prerequisite proof availability.} The named results below
identify specific programme lessons. The
\href{https://kokunoyumeto.github.io/open-math-courses/courses/NT-CFT/proof-dependencies.html\#lesson-12}{prerequisite
record} distinguishes published proofs, supplied owner texts awaiting
publication, and missing full proofs. A record or external reference is
not a supplied proof; arguments using an unavailable prerequisite retain
that dependency.

\subsection{1. Inertia open, Frobenius
discrete}\label{inertia-open-frobenius-discrete}

Let \(G_K=\operatorname{Gal}(K^{\mathrm{sep}}/K)\), and let \(I_K\) be
inertia. The unramified-extension theorem gives \[
1\longrightarrow I_K\longrightarrow G_K
 \xrightarrow{d}\widehat{\mathbf Z}\longrightarrow1.
\tag{1}
\] We choose \(d\) so that an arithmetic Frobenius, acting by
\(x\mapsto x^q\) on the residue field, has degree 1. Define \[
W_K=d^{-1}(\mathbf Z).
\tag{2}
\] Give \(I_K\) its profinite topology and make it open in \(W_K\).
Every coset of \(I_K\) has the topology transported from \(I_K\) by
translation.

\subsubsection{Proposition 12.1. The local Weil
group}\label{proposition-12.1.-the-local-weil-group}

The topology just specified makes \(W_K\) a locally compact, totally
disconnected group. Its quotient by \(I_K\) is the discrete group
\(\mathbf Z\). The inclusion \(W_K\hookrightarrow G_K\) is continuous
and has dense image. The Weil topology is strictly finer than the
topology induced by \(G_K\).

\textbf{Proof.} Choose \(\phi\in G_K\) with \(d(\phi)=1\). Integer
powers of any element of a profinite group extend continuously to
profinite powers: in each finite quotient reduce the exponent modulo
that element's finite order, and take the compatible inverse limit. Thus
\(t\mapsto\phi^t\), \(t\in\widehat{\mathbf Z}\), is a continuous
homomorphism and \(d(\phi^t)=t\). Every element of \(G_K\) has a unique
expression \(i\phi^t\). The corresponding bijection \[
I_K\rtimes\widehat{\mathbf Z}\longrightarrow G_K
\] is a homeomorphism, since its domain is compact and its target
Hausdorff.

On \(W_K\) this expression has \(t\in\mathbf Z\). The prescribed
topology is precisely the product topology of \(I_K\rtimes\mathbf Z\),
with \(\mathbf Z\) discrete. The action of each \(\phi^m\) on \(I_K\) is
continuous, so multiplication and inversion are continuous. The open
subgroup \(I_K\) is compact and totally disconnected, proving the
asserted local properties. The inclusion is continuous on every open
coset. Since \(\mathbf Z\) is dense in \(\widehat{\mathbf Z}\), the
product description proves density.

In the subspace topology the set \(\{\phi^m:m\in\mathbf Z\}\) has the
profinite, rather than discrete, topology. Indeed, every open
neighborhood of 1 in \(G_K\) contains a nonidentity \(\phi^m\): pass to
an open normal subgroup and take a positive multiple of the order of
\(\phi\) in its finite quotient. Such a power has degree \(m\ne0\), so
it is outside inertia. Consequently \(I_K\) is not open in the subspace
topology. ∎

The choice of \(\phi\) splits the group but is not part of its
definition. A different lift changes the splitting, not \(W_K\) or its
topology.

\subsection{2. A topological reciprocity
isomorphism}\label{a-topological-reciprocity-isomorphism}

For any topological group \(H\), write \[
H^{\mathrm{ab}}=H/\overline{[H,H]}.
\] The closure is essential: this is the maximal \textbf{Hausdorff}
abelian quotient.

\subsubsection{Theorem 12.2. Reciprocity for the Weil
group}\label{theorem-12.2.-reciprocity-for-the-weil-group}

Arithmetic local reciprocity induces a topological isomorphism \[
\operatorname{rec}_K:K^\times\xrightarrow{\ \sim\ }W_K^{\mathrm{ab}}.
\tag{3}
\] It carries \(\mathcal O_K^\times\) onto the image of \(I_K\), and a
uniformizer onto a class of arithmetic degree 1.

\textbf{Proof.} Put \(D=\overline{[G_K,G_K]}\). Since \(d\) has abelian
target, \(D\subset I_K\). Density of \(W_K\) implies \[
\overline{[W_K,W_K]}^{\,G_K}=D:
\tag{4}
\] approximate each of the finitely many entries in a product of
commutators by elements of \(W_K\), using continuity of the commutator
map. Both commutator subgroups lie in \(I_K\), where the two topologies
agree. Thus their closure in \(W_K\) is also \(D\). We obtain an
injective map \(W_K/D\to G_K/D\), whose image is exactly the inverse
image of \(\mathbf Z\) under the induced degree map.

Theorem 9.4 identifies \(G_K^{\mathrm{ab}}\) with the completion of
\(K^\times=\pi^{\mathbf Z}\times\mathcal O_K^\times\) for its open
subgroups of finite index: \[
G_K^{\mathrm{ab}}\simeq\widehat{\mathbf Z}\times\mathcal O_K^\times.
\tag{5}
\] The degree is the first coordinate, and reciprocity embeds
\(K^\times\) as \(\mathbf Z\times\mathcal O_K^\times\). By (4), this is
the underlying group \(W_K^{\mathrm{ab}}\).

It remains to check the topology. The subgroup \(I_K/D\) is compact and
open in \(W_K/D\). Its map onto the degree-zero subgroup in (5) is a
continuous bijection from a compact space to a Hausdorff space, hence a
homeomorphism. Reciprocity is therefore a homeomorphism on the unit
group. Each integer-degree coset in both groups is open and is a
translate of that unit group. Translating the homeomorphism on each
coset proves (3). ∎

This argument does not put the subspace topology from
\(G_K^{\mathrm{ab}}\) on \(W_K^{\mathrm{ab}}\). With that topology its
integer coordinate would again fail to be discrete.

\subsection{3. Characters, conductors and Euler
factors}\label{characters-conductors-and-euler-factors}

A character below means a continuous homomorphism to
\(\mathbf C^\times\); it need not be unitary or have finite image.

\subsubsection{Proposition 12.3. Local reciprocity in dimension
one}\label{proposition-12.3.-local-reciprocity-in-dimension-one}

Composition with (3) gives a bijection \[
\{\text{characters of }W_K\}
\longleftrightarrow
\{\text{characters of }K^\times\},\qquad
\chi(a)=\rho(\operatorname{rec}_K(a)).
\tag{6}
\] The restriction of \(\rho\) to inertia has finite image. Unramified
characters correspond to characters trivial on \(\mathcal O_K^\times\).
Their conductors and their arithmetic Euler factors agree.

More precisely, put \(U^0=\mathcal O_K^\times\) and
\(U^j=1+\mathfrak p_K^j\) for \(j\geq1\). The conductor is \[
a(\chi)=\min\{j\geq0:\chi(U^j)=1\}
       =a(\rho).
\tag{7}
\] For \(q^{-s}=\exp(-s\log q)\), define \[
L_{\mathrm{ar}}(s,\rho)
=\det(1-q^{-s}\rho(\phi)\mid V^{I_K})^{-1}.
\] In dimension one this equals \[
L(s,\chi)=
\begin{cases}
(1-\chi(\pi)q^{-s})^{-1},&a(\chi)=0,\\
1,&a(\chi)>0.
\end{cases}
\tag{8}
\]

\textbf{Proof.} A character kills commutators and their closure, because
its target is Hausdorff. It therefore factors through
\(W_K^{\mathrm{ab}}\), and Theorem 12.2 proves the bijection.

The group \(\mathbf C^\times\) has a neighborhood of 1 containing no
nontrivial subgroup. One may use an annulus with argument restricted to
a small arc: powers of an element of nonunit modulus leave the annulus;
powers of a nontrivial unit complex number leave the arc. Continuity
gives an open subgroup \(J\) of the profinite group \(I_K\) whose image
is inside that neighborhood. Its image is a subgroup, hence trivial.
Thus \(\rho(I_K)\) is finite.

To compare conductors without assuming that \(\chi\) has finite image,
write \[
\chi(\pi^m u)=\chi(\pi)^m\chi_0(u),\qquad \chi_0(\pi)=1.
\] The unit character has finite image, so \(\chi_0\) is finite order.
By the existence theorem it factors through a finite abelian extension.
The unit-filtration theorem of lesson 10 identifies the images of
\(U^j\) with its upper ramification groups at integer \(j\). This proves
(7), with conductor zero exactly in the unramified case. The unramified
factor \(\chi(\pi)^m\) does not affect inertia or the conductor.

If inertia acts nontrivially in dimension one, \(V^{I_K}=0\) and the
determinant is 1. Otherwise \(V^{I_K}=V\), and all degree-one lifts give
the scalar \(\chi(\pi)\): the units are killed. This proves (8). ∎

For every \(\lambda\in\mathbf C^\times\), \[
\chi_\lambda(\pi^m u)=\lambda^m
\] is an unramified character. Its Weil representation factors through
the discrete quotient \(W_K/I_K=\mathbf Z\). When \(\lambda=q^{-1}\), it
cannot extend continuously to a character of the compact group \(G_K\):
its image contains the unbounded sequence \(q^m\), \(m\geq0\).

\subsection{4. Rank-one local constants in either
characteristic}\label{rank-one-local-constants-in-either-characteristic}

We record the analytic data as well as the representation. Choose a
nontrivial continuous additive character \(\psi:K\to\mathbf T\) and a
positive additive Haar measure \(dx\). Write \[
N=\max\{r\in\mathbf Z:\psi(\pi^{-r}\mathcal O_K)=1\},
\qquad m=\operatorname{vol}_{dx}(\mathcal O_K).
\tag{9}
\] Continuity and nontriviality make \(N\) a finite integer. Normalize
\(d^\times x\) to give the units volume 1.

Here is the Fourier foundation for all nonarchimedean local fields,
including positive characteristic. Such a \(\psi\) exists: take a
nontrivial additive character of the finite residue field on
\(\mathcal O_K\), and extend successively to \(\pi^{-r}\mathcal O_K\).
At each step the quotient is finite abelian. A circle-valued character
extends across a finite cyclic enlargement by choosing a root of the
prescribed value of a multiple of its generator; adjoining finitely many
generators gives the required extension. The union is \(K\), and the
kernel still contains \(\pi\mathcal O_K\), so the resulting character is
continuous.

Additive Haar measure can be constructed by assigning mass \(m q^{-r}\)
to each coset of \(\pi^r\mathcal O_K\). These assignments are consistent
under subdivision into \(q\) smaller cosets, define the measure on
compact open sets, and then on their Borel completion. Translation
invariance and the scaling identity \(d(ax)=|a|\,dx\) follow on these
cosets and hence on integrable functions.

Let \(\mathcal S(K)\) be the locally constant compactly supported
functions, and use the positive kernel
\(\widehat f(y)=\int f(x)\psi(xy)\,dx\). The annihilator of
\(\pi^r\mathcal O_K\) is \(\pi^{-r-N}\mathcal O_K\): multiplying this
ideal by \(y\) is an ideal on which \(\psi\) is trivial exactly when
\(v(y)+r\geq-N\). Translation and character orthogonality give \[
\widehat{\mathbf1_{a+\pi^r\mathcal O_K}}(y)
=m q^{-r}\psi(ay)\mathbf1_{\pi^{-r-N}\mathcal O_K}(y).
\tag{10}
\] Applying this formula twice gives \[
\widehat{\widehat f}(x)=m^2q^N f(-x).
\tag{11}
\] Every Schwartz function is a finite linear combination of these
indicators, so (10) proves preservation of \(\mathcal S(K)\) and (11)
for every \(f\). The self-dual measure is characterized by
\(m=q^{-N/2}\).

For a multiplicative character \(c\), put \(\check c=c^{-1}|\cdot|\),
and set \[
Z(f,c)=\int_{K^\times}f(x)c(x)\,d^\times x.
\] The integral converges when \(|c(\pi)|<1\). Splitting into shells
\(\pi^j\mathcal O_K^\times\) gives a finite Laurent polynomial plus a
tail. Local constancy makes the tail \(f(0)c(\pi)^J/(1-c(\pi))\) in the
unramified case; unit-character orthogonality makes it zero in the
ramified case. This gives rational continuation in \(c(\pi)\).

The functional equation follows by the same two-variable calculation in
every characteristic. In the strip \(q^{-1}<|c(\pi)|<1\), both \(c\) and
\(\check c\) have absolutely convergent zeta integrals. With
\(d^\times x=\kappa\,dx/|x|\), the substitution \(y=xz\) transforms \[
Z(f,c)Z(\widehat g,\check c)
=\kappa^2\int_{K^\times}c(z)^{-1}
       \left(\int_K f(x)\widehat g(zx)\,dx\right)dz.
\tag{12}
\] The absolute-value double integral is finite by reversing that
substitution. For each fixed \(z\), Fubini on the compact supports gives
\[
\int f(x)\widehat g(zx)\,dx=\int\widehat f(zt)g(t)\,dt.
\] The resulting double integral is also absolutely convergent, as the
reverse substitution now gives the product for \(g,\widehat f\). Hence
(12) equals \(Z(g,c)Z(\widehat f,\check c)\). A test
\(g=c|_{\mathcal O_K^\times}^{-1}\mathbf1_{\mathcal O_K^\times}\) has
zeta integral 1, so this proves \[
Z(\widehat f,\check c)
=\epsilon(c,\psi,dx)\frac{L(\check c)}{L(c)}Z(f,c).
\tag{13}
\] Here \(L(c)=(1-c(\pi))^{-1}\) when \(c\) is unramified and is 1
otherwise. Rational continuation extends (13) beyond the strip. Formula
(11), applied twice, shows that its factor is not identically zero.

For completeness, its finite calculation is \[
\epsilon(c,\psi,dx)=
\begin{cases}
m q^N c(\pi)^N,&a(c)=0,\\
m q^N c(b)\displaystyle\sum_{u\in(\mathcal O_K/\mathfrak p_K^a)^\times}
 c(u)^{-1}\psi(u/b),&a(c)=a\geq1,
\end{cases}
\quad v(b)=a+N.
\tag{14}
\] In the first case use \(f=\mathbf1_{\mathcal O_K}\) and (10). In the
second use the preceding unit-character test. Translation by
\(\mathfrak p_K^a\) forces its transform to be supported on
\(\pi^{-a-N}\mathcal O_K\). It vanishes at every strictly higher
valuation: for \(a\geq2\) multiply the unit variable by \(v\in U^{a-1}\)
with \(c(v)\ne1\), which leaves the additive factor unchanged; for
\(a=1\), ordinary orthogonality on the units suffices. At \(y=w/b\) the
transform is \[
m q^{-a}c(w)\sum_u c(u)^{-1}\psi(u/b).
\] Its dual zeta integral on that one shell multiplies this by
\(c(b)|b|^{-1}c(w)^{-1}\), proving (14). Replacing \(b\) by \(vb\),
\(v\) a unit, multiplies its two factors by \(c(v)\) and \(c(v)^{-1}\),
respectively.

The Gauss sum is nonzero. Indeed, applying (11) again to this test gives
the product of the sums for the unit character and its inverse equal to
\(c(-1)q^a\); complex conjugation gives absolute value \(q^{a/2}\) for
either sum. This proves nonvanishing in (14) for every \(c(\pi)\ne0\),
not just meromorphically. All steps used finite quotients, compact open
cosets and (10); none required characteristic zero. They also explain
exactly how the finite-place Tate formulas extend to function-field
completions.

For the corresponding Weil character define \[
\epsilon_{\mathrm{ar}}(s,\rho,\psi,dx)
=\epsilon(\chi|\cdot|^s,\psi,dx).
\tag{15}
\] This transports the full functional equation, rather than just an
unspecified scalar. In particular, with the measure fixed, \[
\epsilon(c,\psi,t\,dx)=t\epsilon(c,\psi,dx),\qquad
\epsilon(c,\psi_r,dx)=c(r)|r|^{-1}\epsilon(c,\psi,dx).
\tag{16}
\] These follow by scaling the Fourier transform and by substituting
\(y\mapsto ry\) in its zeta integral. For self-dual measure, (11) and
cancellation of the two \(L\)-ratios give \[
\epsilon(c,\psi,dx_\psi)\epsilon(c^{-1}|\cdot|,\psi,dx_\psi)=c(-1).
\tag{17}
\] Thus duality involves an inverse character \textbf{and} a norm twist;
it is not a rule that simply inverts epsilon.

The higher-dimensional extension is Deligne's Theorem 4.1: there is a
unique family of nonzero constants for finite-dimensional smooth complex
Weil representations, multiplicative in exact sequences, scaled by
\(t^{\dim V}\) under \(dx\mapsto t\,dx\), compatible with induction of
virtual representations of dimension zero through finite separable
extensions using \(\psi\circ\operatorname{Tr}\), and equal to Tate's
constants in dimension one under geometric reciprocity. Its
representation-theoretic proof belongs to
\href{https://kokunoyumeto.github.io/open-math-courses/courses/LG-GAL/local-l-factors-and-epsilon-factors.html\#3-existence-and-uniqueness-of-local-constants}{Local
L-factors and epsilon-factors of Weil group representations}, section 3.
Theorem 3.0 and sections 3A–3D supply the existence argument, including the Brauer-relation check and arbitrary smooth representations; Theorem 3.2 supplies uniqueness. Our proofs of (3), (6)--(8) and
(13)--(17) do not use that higher-dimensional existence assertion.

\subsection{5. The real and complex Weil
groups}\label{the-real-and-complex-weil-groups}

\subsubsection{Proposition 12.4. Archimedean
abelianizations}\label{proposition-12.4.-archimedean-abelianizations}

Set \(W_{\mathbf C}=\mathbf C^\times\) with its usual topology. The real
Weil group has two components \[
W_{\mathbf R}=\mathbf C^\times\sqcup j\mathbf C^\times,\qquad
j^2=-1,\qquad jzj^{-1}=\bar z.
\tag{20}
\] It is a locally compact group, and \[
W_{\mathbf C}^{\mathrm{ab}}=\mathbf C^\times,\qquad
W_{\mathbf R}^{\mathrm{ab}}\simeq\mathbf R^\times,
\quad z\longmapsto|z|^2,\quad j\longmapsto-1
\tag{21}
\] are topological isomorphisms.

\textbf{Proof.} Realize \(z\) as the complex matrix
\(\operatorname{diag}(z,\bar z)\) and \(j\) as
\(\left(\begin{smallmatrix}0&-1\\1&0\end{smallmatrix}\right)\). Their
two sets of matrices are closed under multiplication and satisfy (20),
so associativity follows from matrix multiplication. Each component has
the usual topology of \(\mathbf C^\times\), and the relations show
continuity of the group operations.

The indicated map \(p:W_{\mathbf R}\to\mathbf R^\times\) is a surjective
homomorphism: the relation \(j^2=-1\) maps to \((-1)^2=|-1|^2=1\), and
conjugation preserves \(|z|^2\). Its kernel is the circle \(S^1\) in the
first component. The commutators \([j,z]=\bar z/z\) fill that circle,
while every commutator lies in the kernel because the quotient is
abelian. Thus the commutator subgroup equals \(S^1\), already closed.
The inverse map on quotient classes is represented by \(\sqrt{x}\) for
\(x>0\) and \(j\sqrt{-x}\) for \(x<0\); these representatives vary
continuously on the two open components. This proves the real
topological assertion. The complex assertion is immediate because
\(\mathbf C^\times\) is abelian. ∎

The real sign character is 1 on \(\mathbf C^\times\) and is \(-1\) on
\(j\mathbf C^\times\). The map \(z\mapsto|z|^2\) in (21) also displays
compatibility with the norm \(\mathbf C^\times\to\mathbf R^\times\).

All real multiplicative characters have the form
\(\operatorname{sgn}(x)^e|x|^t\), \(e\in\{0,1\}\), \(t\in\mathbf C\).
All complex ones have the form \((z/|z|)^\ell|z|_{\mathbf C}^t\),
\(\ell\in\mathbf Z\), with \(|z|_{\mathbf C}=|z|^2\). To see this, a
continuous character of \(\mathbf R_{>0}\) becomes a character of
\((\mathbf R,+)\) under logarithm; lifting near 1 and subdividing
intervals gives \(x\mapsto e^{ax}\). A character of \(S^1\) lifted along
\(e^{i\theta}\) has form \(e^{ib\theta}\), and periodicity forces
\(b\in\mathbf Z\). A character of \(\{\pm1\}\) has values \(1,(-1)^e\).

Transporting the proved archimedean Tate factors through (21), for the
negative additive characters \(\psi_{\mathbf R}(x)=e^{-2\pi ix}\),
\(\psi_{\mathbf C}(z)=e^{-4\pi i\operatorname{Re}z}\), and self-dual
measures \(dx\), \(2\,d(\operatorname{Re}z)\,d(\operatorname{Im}z)\),
respectively, gives \[
\begin{array}{c|c|c}
\chi&L(s,\chi)&\epsilon(s,\chi)\\ \hline
\operatorname{sgn}^e|\cdot|^t&
\Gamma_{\mathbf R}(s+t+e)&(-i)^e\\
(z/|z|)^\ell|\cdot|_{\mathbf C}^t&
\Gamma_{\mathbf C}(s+t+|\ell|/2)&(-i)^{|\ell|}
\end{array}
\tag{22}
\] where \(\Gamma_{\mathbf R}(u)=\pi^{-u/2}\Gamma(u/2)\) and
\(\Gamma_{\mathbf C}(u)=2(2\pi)^{-u}\Gamma(u)\). The Gaussian and
polynomial Fourier proofs are \emph{Tate's local theory at the infinite
places}, Propositions 8.2--8.3 and Theorem 8.1. Replacing these negative
additive characters by their positives multiplies epsilon by
\(\chi(-1)\), giving \(i^e\) and \(i^{|\ell|}\). This phase change is
independent of a Frobenius convention.

\subsection{6. Translating Frobenius
conventions}\label{translating-frobenius-conventions}

\subsubsection{Proposition 12.5. Arithmetic and geometric
reciprocity}\label{proposition-12.5.-arithmetic-and-geometric-reciprocity}

For nonarchimedean \(K\), geometric reciprocity is \[
\operatorname{rec}^{\mathrm{geo}}_K(a)=\operatorname{rec}_K(a)^{-1}.
\tag{18}
\] It sends a uniformizer to a geometric Frobenius class. For a fixed
Weil character \(\rho\), its multiplicative character under (18) is
\(\chi^{-1}\). For a fixed multiplicative character \(\chi\), the
corresponding Weil characters in the two conventions are dual.
Conductors are unchanged. For any complex Weil representation with
finite inertia image, \[
L_{\mathrm{geo}}(s,V)=L_{\mathrm{ar}}(s,V^\vee),
\tag{19}
\] where each side uses the indicated Frobenius on the same underlying
representation or its dual.

\textbf{Proof.} Inversion is a continuous automorphism of the abelian
group \(W_K^{\mathrm{ab}}\), and it reverses degree. Evaluating \(\rho\)
on (18) proves the character statement. Inverse characters have the same
kernel on each unit group, proving the conductor statement.

For (19), averaging over the finite inertia image identifies
\((V^\vee)^{I_K}\) with the dual of \(V^{I_K}\): the averaging
idempotent has image \(V^{I_K}\), and its dual has image the dual of
that summand. Frobenius preserves both summands. Its operator on the
dual invariants is the transpose inverse of its operator on the original
invariants, which has the same determinant polynomial as the inverse
operator. A geometric lift is the inverse of an arithmetic lift up to
inertia, which is trivial on these invariants. This proves (19). ∎

For example an unramified fixed representation with arithmetic
eigenvalue \(\lambda\) has arithmetic factor \((1-\lambda q^{-s})^{-1}\)
and geometric factor \((1-\lambda^{-1}q^{-s})^{-1}\). If one changes the
reciprocity map and also dualizes the representation so that the
multiplicative character is held fixed, its Tate factor stays the same.

Deligne fixes geometric reciprocity in §2.3 and explains its
cohomological Euler-factor convention in §3.6. When using his factors,
evaluate the multiplicative character through his reciprocity map; do
not keep an arithmetic character identification while inserting a
geometric Frobenius. Additive characters and measures are separate
choices, governed by (16). At archimedean places there is no
residue-field Frobenius and hence no arithmetic/geometric choice of this
kind.

\subsection{7. Global Weil groups and the idèle class
group}\label{global-weil-groups-and-the-iduxe8le-class-group}

Let \(C_F=\mathbb A_F^\times/F^\times\) for a global field \(F\). The
local theorem has a global analogue \[
C_F\simeq W_F^{\mathrm{ab}},
\tag{23}
\] compatible with the maps \(F_v^\times\to C_F\) and local Weil groups.

For a function field over its full finite constant field, the group is
again a subgroup of the absolute Galois group: take the inverse image of
the integral powers of constant-field Frobenius and make its geometric
kernel open. The proof of (23) uses global reciprocity, the compact
degree-zero subgroup of \(C_F\), and its integral degree quotient. The
complete proof is
\href{../global-existence-and-the-idelic-class-field-correspondence.html}{Global
existence and the idèlic class field correspondence}, sections 5--8,
equations (14)--(34). It proves injectivity even for
characteristic-power extensions, identifies the compact degree-zero
kernel, and verifies that the Hausdorff commutator quotient has exactly
the required topology and local compatibility.

A number field has no finite constant field whose Frobenius supplies
such a subgroup. Its global Weil group is instead constructed from
relative extensions \[
1\longrightarrow C_L\longrightarrow W_{L/F}
 \longrightarrow\operatorname{Gal}(L/F)\longrightarrow1
\tag{24}
\] defined by the global class-formation fundamental classes, with
compatible transition maps as \(L/F\) runs through finite Galois
extensions.
\href{../brauer-groups-of-local-and-global-fields.html}{Brauer groups of
local and global fields}, Theorem 24.6 and section 7, proves the
fundamental-class theorem, constructs these extensions and their
coherent quotient maps, and proves local compactness, the Hausdorff
abelianization and local compatibility in (23). Its equations (29)--(39)
give the complete number-field construction. The function-field
construction is the preceding proof in lesson 17. These global arguments
use reciprocity established later in the course; the local proofs above
are independent of them. The classical number-field construction is due
to Weil; Deligne §2.5 identifies this source rather than constructing it
there.

\subsection{8. Exercises with solutions}\label{exercises-with-solutions}

\subsubsection{Exercise 1. Density and the topology ---
easy}\label{exercise-1.-density-and-the-topology-easy}

Prove density of \(W_K\) and openness of inertia. Show why density does
not give the Weil group the subspace topology.

\textbf{Solution.} In the splitting
\(G_K=I_K\rtimes\widehat{\mathbf Z}\), the subset \(W_K\) is
\(I_K\rtimes\mathbf Z\). Every product neighborhood meets it because an
open subset of \(\widehat{\mathbf Z}\) contains an integer. In the Weil
topology \(I_K\) is open by construction. In the Galois subspace
topology every identity neighborhood contains a nonzero integral power
of \(\phi\), by taking a multiple of its order in a finite quotient.
Therefore inertia is not open in that topology.

\subsubsection{Exercise 2. The real abelianization ---
medium}\label{exercise-2.-the-real-abelianization-medium}

Compute the closed commutator subgroup of \(W_{\mathbf R}\) and recover
the sign representation.

\textbf{Solution.} For \(z=re^{i\theta}\), \([j,z]=e^{-2i\theta}\), so
these commutators fill \(S^1\). The homomorphism \(z\mapsto r^2\),
\(j\mapsto-1\) has precisely this kernel; hence all commutators are in
\(S^1\) and the closed commutator subgroup equals \(S^1\). The quotient
is \(\mathbf R^\times\), with the continuous inverse classes described
in Proposition 12.4. Pulling back sign gives 1 on the first component
and \(-1\) on the second; the relation \(j^2=-1\) is respected since
both sides have character value 1.

\subsubsection{Exercise 3. Matching conductors ---
medium}\label{exercise-3.-matching-conductors-medium}

Let \(\rho\) correspond to \(\chi\). Prove equality of their conductors
even when \(\chi(\pi)\) has infinite order.

\textbf{Solution.} The unit restriction has finite image by the
no-small-subgroup argument. Remove the unramified character taking
\(\pi\) to \(\chi(\pi)\), leaving a finite-order character \(\chi_0\)
with the same restriction to every \(U^j\). It factors through a finite
abelian extension by local existence. Reciprocity sends \(U^j\) to its
integer upper ramification group, so the least \(j\) on which \(\chi_0\)
is trivial equals the Weil conductor. Restoring the unramified factor
changes neither side.

\subsubsection{Exercise 4. Proving the homeomorphism ---
hard}\label{exercise-4.-proving-the-homeomorphism-hard}

A continuous bijection need not be a homeomorphism. Supply the
topological step in Theorem 12.2 and contrast the embedding in
\(G_K^{\mathrm{ab}}\).

\textbf{Solution.} Write \(D=\overline{[G_K,G_K]}\subset I_K\). In
\(W_K/D\), the subgroup \(I_K/D\) is compact and open. Reciprocity
restricts to a continuous bijection \(\mathcal O_K^\times\to I_K/D\);
compactness and the Hausdorff property make it a homeomorphism. Both
groups are disjoint unions of open translates indexed by \(\mathbf Z\).
On every translate, multiplication transports that homeomorphism, so the
inverse map is continuous on an open cover and hence everywhere. In
\(G_K^{\mathrm{ab}}\) the degree coordinate instead lies in
\(\widehat{\mathbf Z}\). Its integer subgroup is dense and has a
nondiscrete induced topology, so the original dense embedding is not the
same topological assertion.

\subsection{Editable edition}\label{editable-edition}

The reading edition provides the complete LaTeX source of this lesson,
the cumulative course LaTeX and the editable source ZIP. The archive
contains all twenty-four Markdown lessons, complete LaTeX bodies,
original diagrams, metadata and reproduction instructions.

\subsection{References}\label{references}

The local Weil construction and one-dimensional formulas use the written
rank-one Fourier prerequisites. Global Weil constructions are supplied
in lessons 17 and 24. General higher-dimensional epsilon existence is supplied by LG-GAL-06, Theorem 3.0 and sections 3A–3D; its uniqueness theorem is Theorem 3.2. These are programme proofs, distinct from the external reading sources.

\begin{itemize}
\tightlist
\item
  \href{https://math.mit.edu/~poonen/786/notes.pdf}{Bjorn Poonen, Tate's
  Thesis, MIT 18.786 lecture notes (2015)}.
\item
  \href{https://ncatlab.org/nlab/files/TateNumberTheory.pdf}{John Tate,
  Number theoretic background (1979), freely available paper}.
\item
  \href{https://publications.ias.edu/sites/default/files/Number20.pdf}{Pierre
  Deligne, Les constantes des équations fonctionnelles des fonctions L
  (1973), IAS archive}.
\end{itemize}

The \href{../free-proofs.html}{proof guide} gives the lesson sequence
and the exact prerequisite record. External references accompany the
written arguments.

\end{document}
